\documentclass[11pt, a4paper]{article}

% --- Page Geometry & Layout ---
\usepackage[a4paper, margin=2.5cm, top=2.8cm, bottom=2.8cm]{geometry}

% --- Font & Language Setup ---
\usepackage{fontspec}
\usepackage[english, bidi=basic, provide=*]{babel}
\babelprovide[import, onchar=ids fonts]{portuguese}
\babelfont{rm}{Noto Sans}
\babelfont{sf}{Noto Sans}
\babelfont{tt}{Noto Sans Mono}

% --- Packages ---
\usepackage{amsmath, amssymb, amsthm}
\usepackage{booktabs}
\usepackage{caption}
\usepackage{xcolor}
\usepackage{microtype}
\usepackage{fancyhdr}
\usepackage{titlesec}
\usepackage{enumitem}
\usepackage{listings}
\usepackage[hidelinks, pdfencoding=auto]{hyperref}

% --- Color Palette ---
\definecolor{primaryColor}{RGB}{15, 23, 42}      % Dark Slate
\definecolor{accentColor}{RGB}{2, 132, 199}      % Sky Blue accent
\definecolor{neutralLight}{RGB}{248, 250, 252}   % Off-white / Slate 50
\definecolor{borderColor}{RGB}{226, 232, 240}    % Slate 200
\definecolor{codeBg}{RGB}{241, 245, 249}         % Slate 100

% --- Header & Footer Configuration ---
\pagestyle{fancy}
\fancyhf{}
\renewcommand{\headrulewidth}{0.4pt}
\renewcommand{\footrulewidth}{0.4pt}
\fancyhead[L]{\small\color{primaryColor}\textbf{Criptografia RSA e Chat Seguro}}
\fancyhead[R]{\small\color{gray}Relatório Teórico}
\fancyfoot[L]{\small\color{gray}Disciplina: Segurança da Informação / Criptografia Aplicada}
\fancyfoot[R]{\small\color{primaryColor}\thepage}

% --- Section Styling ---
\titleformat{\section}
  {\normalfont\Large\bfseries\color{primaryColor}}{\thesection.}{0.6em}{}
  [\color{accentColor}\titlerule]
\titleformat{\subsection}
  {\normalfont\large\bfseries\color{primaryColor}}{\thesubsection}{0.6em}{}

% --- Code Block / Log Box Styling ---
\lstdefinestyle{logstyle}{
    backgroundcolor=\color{codeBg},
    basicstyle=\ttfamily\footnotesize\color{primaryColor},
    breakatwhitespace=false,
    breaklines=true,
    captionpos=b,
    keepspaces=true,
    showspaces=false,
    showstringspaces=false,
    showtabs=false,
    tabsize=2,
    frame=single,
    rulecolor=\color{borderColor}
}

% --- Document Metadata ---
\title{
    \vspace{-1.5cm}
    \color{primaryColor}\huge\bfseries Criptografia RSA e\\Implementação de Chat Seguro\\[0.3em]
    \large\color{accentColor}Relatório Teórico --- Trabalho Prático
}
\author{\textbf{Disciplina:} Segurança da Informação / Criptografia Aplicada \\ \textbf{Aluno:} \underline{\hspace{6cm}} \\ \textbf{Data:} Outubro de 2026}
\date{}

\begin{document}

\maketitle
\thispagestyle{fancy}

\vspace{-0.5cm}

\section{Introdução}
Este relatório apresenta a fundamentação teórica do algoritmo RSA, a análise de sua segurança e a descrição da implementação desenvolvida neste trabalho: uma biblioteca RSA escrita do zero em Python (módulo \texttt{rsa\_core.py}), sem uso de bibliotecas de criptografia de terceiros, e um chat cliente-servidor (\texttt{chat\_server.py} e \texttt{chat\_client.py}) em que toda a conversa é cifrada com RSA e a troca de chaves públicas é feita pela rede.

O texto segue o roteiro exigido: história e importância do RSA, explicação matemática passo a passo, processo de cifragem e decifragem com exemplo numérico, e análise de segurança. A base bibliográfica principal é o material indicado pelo professor (arquivo disponibilizado em link do Google Drive), complementado pelas fontes citadas nas referências.

\section{História e importância do RSA}
Até a década de 1970, toda a criptografia prática era simétrica: emissor e receptor precisavam compartilhar previamente a mesma chave secreta, o que cria o chamado problema da distribuição de chaves. Em 1976, Whitfield Diffie e Martin Hellman propuseram publicamente a ideia de criptografia de chave pública, em que cada usuário possui um par de chaves: uma pública, que pode ser divulgada, e uma privada, mantida em segredo. Eles, porém, não apresentaram uma função que realizasse a cifragem assimétrica.

Em 1977, Ron Rivest, Adi Shamir e Leonard Adleman, do MIT, criaram o algoritmo que leva as iniciais de seus sobrenomes e o publicaram em 1978 [1]. O RSA baseia-se na assimetria entre duas tarefas: multiplicar dois números primos grandes é fácil, mas fatorar o produto é, até hoje, computacionalmente inviável quando os primos são suficientemente grandes. Posteriormente tornou-se público que Clifford Cocks, no serviço secreto britânico GCHQ, já havia descoberto um esquema equivalente em 1973, mantido sob sigilo até 1997.

O RSA foi o primeiro algoritmo amplamente adotado que permite cifrar e também assinar digitalmente. Sua importância na criptografia moderna está em:
\begin{enumerate}[label=(\roman*)]
    \item resolver a distribuição de chaves, pois a chave pública pode trafegar em canais abertos;
    \item viabilizar assinaturas digitais e certificados, base das infraestruturas de chave pública (PKI) e do protocolo TLS/HTTPS; e
    \item servir de modelo para o desenho de outros sistemas assimétricos.
\end{enumerate}
Na prática, por ser lento, o RSA costuma ser usado para proteger uma pequena quantidade de dados, como a chave simétrica de uma sessão, enquanto um algoritmo simétrico (por exemplo, AES) cifra o volume principal da comunicação.

\section{Fundamentação matemática passo a passo}

\subsection{Visão geral}
O RSA trabalha com aritmética modular sobre números inteiros. A segurança vem da dificuldade de fatorar o módulo $n$; a correção vem de um resultado clássico da teoria dos números, o Teorema de Euler. A geração das chaves segue cinco passos.

\subsection{Passo 1 --- Geração de dois números primos}
Escolhem-se dois primos distintos, $p$ e $q$, grandes e aleatórios (hoje, cada um com 1024 bits ou mais). Como não existe fórmula para produzir primos, o método é sortear um candidato ímpar com o tamanho desejado e testá-lo. No trabalho, o teste usado é o de Miller-Rabin, baseado no Pequeno Teorema de Fermat: se $n$ é primo e $1 < a < n$, então $a^{n-1} \pmod n = 1$. Escrevendo $n-1 = 2^s \cdot d$ ($d$ ímpar), o teste verifica se $a^d \pmod n$ vale $1$ ou $n-1$, ou se algum dos quadrados sucessivos vale $n-1$. Se nenhuma dessas condições ocorre, $n$ é composto com certeza. Cada rodada com base aleatória $a$ detecta um composto com probabilidade de pelo menos $3/4$; com 40 rodadas a chance de erro é inferior a $4^{-40}$. Pelo Teorema dos Números Primos, a densidade de primos próximos de $N$ é cerca de $1/\ln(N)$, de modo que poucas centenas de sorteios bastam.

\subsection{Passo 2 --- Cálculo do módulo $n$}
O módulo é o produto dos primos:
\begin{equation}
    n = p \cdot q
\end{equation}
O número $n$ faz parte das duas chaves e é público. O tamanho em bits de $n$ é o ``tamanho da chave RSA'' (por exemplo, 2048 bits).

\subsection{Passo 3 --- Função totiente de Euler, $\phi(n)$}
A função totiente $\phi(n)$ conta quantos inteiros entre 1 e $n$ são coprimos com $n$. Para um número primo $p$, $\phi(p) = p - 1$, e a função é multiplicativa para números coprimos. Logo:
\begin{equation}
    \phi(n) = (p - 1) \cdot (q - 1)
\end{equation}
Calcular $\phi(n)$ é trivial para quem conhece $p$ e $q$, mas é equivalente a fatorar $n$ para quem não os conhece. Por isso $p$, $q$ e $\phi(n)$ são descartados ou mantidos em segredo.

\subsection{Passo 4 --- Escolha do expoente público $e$}
Escolhe-se $e$ tal que $1 < e < \phi(n)$ e $\text{mdc}(e, \phi(n)) = 1$, isto é, $e$ é coprimo com $\phi(n)$. A condição é necessária para que exista o inverso modular de $e$. O MDC é calculado pelo algoritmo de Euclides: $\text{mdc}(a,b) = \text{mdc}(b, a \bmod b)$. Na prática adota-se $e = 65537$ ($2^{16} + 1$), um primo grande o bastante para ser seguro e com apenas dois bits 1, o que torna a cifragem rápida.

\subsection{Passo 5 --- Cálculo do expoente privado $d$}
O expoente privado é o inverso multiplicativo de $e$ módulo $\phi(n)$:
\begin{equation}
    d \cdot e \equiv 1 \pmod{\phi(n)}
\end{equation}
Ele é obtido pelo algoritmo de Euclides estendido, que encontra $x$ e $y$ tais que $e \cdot x + \phi(n) \cdot y = \text{mdc}(e, \phi(n)) = 1$ (identidade de Bézout). Reduzindo módulo $\phi(n)$, obtém-se $e \cdot x \equiv 1$, e portanto $d = x \bmod \phi(n)$.

Ao final, a chave pública é o par $(e, n)$ e a chave privada é o par $(d, n)$.

\subsection{Por que a decifragem funciona}
Como $e \cdot d \equiv 1 \pmod{\phi(n)}$, existe um inteiro $k$ com $e \cdot d = 1 + k \cdot \phi(n)$. Para uma mensagem $m$ coprima com $n$, o Teorema de Euler garante $m^{\phi(n)} \equiv 1 \pmod n$. Assim:
\begin{equation}
    (m^e)^d = m^{e \cdot d} = m^{1 + k \cdot \phi(n)} = m \cdot (m^{\phi(n)})^k \equiv m \cdot 1^k \equiv m \pmod n
\end{equation}
O resultado também vale quando $m$ não é coprimo com $n$ (caso raríssimo para primos grandes), o que se demonstra aplicando o Teorema de Fermat separadamente módulo $p$ e módulo $q$ e combinando os resultados pelo Teorema Chinês do Resto. Portanto, elevar a $e$ e depois a $d$ devolve a mensagem original.

\section{Criptografia e descriptografia}
Com a chave pública $(e, n)$ do destinatário, o remetente transforma a mensagem em um inteiro $m$, com $0 \le m < n$, e calcula:
\begin{equation}
    c = m^e \bmod n \quad \text{(cifragem)}
\end{equation}
O destinatário usa sua chave privada $(d, n)$ para recuperar:
\begin{equation}
    m = c^d \bmod n \quad \text{(decifragem)}
\end{equation}
Como os expoentes são enormes, usa-se a exponenciação modular rápida (\textit{square-and-multiply}): percorrem-se os bits do expoente, elevando a base ao quadrado a cada passo e multiplicando pelo resultado parcial quando o bit é 1, sempre reduzindo módulo $n$. O custo é proporcional a $\log(\text{expoente})$, e não ao valor do expoente.

\subsection{Exemplo numérico com números pequenos}
Usaremos os valores didáticos clássicos (números pequenos só para permitir conferir à mão):

\begin{table}[h!]
    \centering
    \caption{Passos da geração de chaves no exemplo numérico}
    \begin{tabular}{lll}
        \toprule
        \textbf{Passo} & \textbf{Cálculo} & \textbf{Resultado} \\
        \midrule
        Primos & $p=61$, $q=53$ & $p, q$ \\
        Módulo & $n = 61 \times 53$ & $n = 3233$ \\
        Totiente & $\phi(n) = 60 \times 52$ & $\phi(n) = 3120$ \\
        Expoente público & $e = 17$, $\text{mdc}(17, 3120) = 1$ & $e = 17$ \\
        Expoente privado & $d = 17^{-1} \bmod 3120$ & $d = 2753$ \\
        Chave pública & $(e, n)$ & $(17, 3233)$ \\
        Chave privada & $(d, n)$ & $(2753, 3233)$ \\
        \bottomrule
    \end{tabular}
\end{table}

\noindent \textbf{Verificação do inverso:} $17 \times 2753 = 46801 = 15 \cdot 3120 + 1$, logo $2753 \equiv 1 \pmod{3120}$. Correto.\\[0.3em]
\noindent \textbf{Cifrando $m=65$:} $c = 65^{17} \bmod 3233 = 2790$.\\[0.3em]
\noindent \textbf{Decifrando $c=2790$:} $m = 2790^{2753} \bmod 3233 = 65$.\\[0.3em]
A mensagem original é recuperada. Outro exemplo: $m=42$ é cifrado como $c=2557$ e decifrado de volta para 42 (valores conferidos pelo código \texttt{rsa\_core.py}).

\subsection{Como converter texto em números}
O RSA cifra inteiros menores que $n$, e não texto. Na implementação, a string é convertida para bytes UTF-8 (o que suporta acentos e emojis) e dividida em blocos de $(k-2)$ bytes, em que $k$ é o tamanho de $n$ em bytes. Cada bloco recebe o byte marcador \texttt{0x01} à frente, o que preserva bytes nulos iniciais e garante $m < n$. Cada bloco vira um inteiro, é cifrado separadamente e o resultado é escrito em hexadecimal com largura fixa para permitir a separação dos blocos na decifragem.

\section{Análise de segurança}

\subsection{O problema da fatoração}
Quem conhece apenas a chave pública $(e, n)$ pode, em tese, descobrir a chave privada se conseguir fatorar $n$ em $p$ e $q$: calcula-se $\phi(n)$ e depois $d$ pelo Euclides estendido. A segurança prática do RSA depende, portanto, da inviabilidade de fatorar $n$. Fatorar por divisão tentativa exige cerca de $\sqrt{n}$ operações, ou seja, $2^{b/2}$ para um módulo de $b$ bits. A tabela abaixo mostra tempos medidos com um fatorador ingênuo (divisão tentativa) em um computador comum, durante os testes deste trabalho: a cada 4 bits adicionais o tempo quadruplica.

\begin{table}[h!]
    \centering
    \caption{Desempenho da fatoração por divisão tentativa segundo o tamanho do módulo}
    \begin{tabular}{cccccc}
        \toprule
        \textbf{Módulo $n$ (bits)} & 24 & 32 & 40 & 44 & 48 \\
        \midrule
        \textbf{Tempo (s)} & 0,0002 & 0,0027 & 0,039 & 0,158 & 0,501 \\
        \bottomrule
    \end{tabular}
\end{table}

Extrapolando: um módulo de 1024 bits exigiria da ordem de $2^{512}$ operações por divisão tentativa, número astronômico. Algoritmos melhores existem, em especial o crivo geral de corpos de números (GNFS), de custo sub-exponencial, mas mesmo ele é inviável para tamanhos atuais de chave.

\subsection{Tamanho das chaves e marcos de fatoração}
Os desafios públicos de fatoração mostram a evolução. O módulo RSA-768 (768 bits, 232 dígitos decimais) foi fatorado em 2009, com esforço de milhares de anos-núcleo de CPU. O RSA-250 (829 bits) foi fatorado em 2020 [3]. Esses recordes mostram que chaves de 512 bits ou menos são completamente inseguras e que 1024 bits já não é considerado adequado para uso real.

\begin{table}[h!]
    \centering
    \caption{Níveis de segurança por tamanho de chave segundo o NIST}
    \begin{tabular}{lll}
        \toprule
        \textbf{Tamanho de $n$} & \textbf{Nível de segurança aproximado} & \textbf{Situação} \\
        \midrule
        512 bits & baixo (fatorável com poucos recursos) & Inseguro \\
        1024 bits & $\sim 80$ bits & Desaconselhado; usado só nos testes deste trabalho \\
        2048 bits & $\sim 112$ bits & Mínimo recomendado atualmente \\
        3072 bits & $\sim 128$ bits & Recomendado para proteção de longo prazo \\
        \bottomrule
    \end{tabular}
\end{table}

A correspondência entre tamanho de chave e nível de segurança segue as recomendações do NIST [4]. Dobrar o tamanho da chave aumenta muito o custo de quebra, mas também torna mais lenta a geração de chaves e a decifragem.

\subsection{Outros pontos de atenção}
\begin{itemize}[leftmargin=*]
    \item \textbf{Computação quântica.} O algoritmo de Shor [5] fatora inteiros em tempo polinomial em um computador quântico suficientemente grande, o que tornaria o RSA inseguro. Por isso há migração para algoritmos pós-quânticos padronizados pelo NIST, e as diretrizes do órgão preveem a descontinuação gradual do RSA.
    \item \textbf{Escolha ruim dos primos.} Primos pequenos, primos próximos entre si (vulnerável ao método de Fermat), ou primos gerados com aleatoriedade fraca permitem recuperar a chave privada. Por isso o código usa o módulo \texttt{secrets} e força os dois bits mais altos de cada primo.
    \item \textbf{RSA sem preenchimento (\textit{textbook RSA}).} A versão simples implementada aqui é determinística (a mesma mensagem gera sempre a mesma cifra, permitindo ataques de dicionário) e maleável (é possível manipular a cifra de modo previsível). Em sistemas reais usa-se preenchimento aleatório, como o OAEP do padrão PKCS\#1 [6], além de resumo e assinatura para integridade. Outros ataques conhecidos estão resumidos em [7].
    \item \textbf{Ausência de autenticação das chaves.} No chat deste trabalho um atacante ativo, que interceptasse o handshake e trocasse as chaves públicas pelas suas (ataque \textit{man-in-the-middle}), poderia ler a conversa. A solução real é autenticar chaves públicas por meio de certificados digitais.
\end{itemize}

\section{Implementação}

\subsection{Módulo \texttt{rsa\_core.py} (Parte 2)}
O módulo implementa, sem bibliotecas de criptografia, as funções abaixo (a única importação é \texttt{secrets}, da biblioteca padrão, usada somente como fonte de bytes aleatórios):

\begin{table}[h!]
    \centering
    \caption{Funções desenvolvidas no módulo \texttt{rsa\_core.py}}
    \begin{tabular}{ll}
        \toprule
        \textbf{Função} & \textbf{Papel} \\
        \midrule
        \texttt{mdc(a, b)} & Máximo divisor comum (Euclides) \\
        \texttt{mdc\_estendido(a, b)} & Coeficientes de Bézout, base do inverso modular \\
        \texttt{inverso\_modular(a, m)} & Calcula $d = e^{-1} \bmod \phi(n)$ \\
        \texttt{exp\_modular(b, e, m)} & Exponenciação modular rápida \\
        \texttt{eh\_primo(n)} & Teste de primalidade (Miller-Rabin, 40 rodadas) \\
        \texttt{gerar\_primo(bits)} & Primo aleatório com o tamanho pedido \\
        \texttt{gerar\_chaves(bits)} & Retorna $((e, n), (d, n))$ gerados dinamicamente \\
        \texttt{criptografar(texto, pub)} & $\text{str} \rightarrow \text{UTF-8} \rightarrow \text{blocos} \rightarrow c = m^e \bmod n \rightarrow \text{hex}$ \\
        \texttt{descriptografar(cifra, priv)} & $\text{hex} \rightarrow m = c^d \bmod n \rightarrow \text{blocos} \rightarrow \text{str}$ \\
        \bottomrule
    \end{tabular}
\end{table}

\subsection{Chat criptografado (Parte 3)}
O chat usa sockets TCP. Os arquivos \texttt{chat\_server.py} e \texttt{chat\_client.py} apenas estabelecem a conexão; depois dela, ambos executam a mesma rotina (\texttt{chat\_common.py}), de modo que a conversa é simétrica. O fluxo é:
\begin{enumerate}
    \item \textbf{Geração de chaves:} cada nó gera seu par RSA de 1024 bits ao iniciar.
    \item \textbf{Handshake:} cada nó envia sua chave pública na linha \texttt{PUBKEY e n} e lê a do outro.
    \item \textbf{Envio:} o texto digitado é cifrado com a chave pública do destinatário e enviado como \texttt{MSG cifra\_hex}.
    \item \textbf{Recebimento:} uma thread separada recebe a cifra e a decifra com a chave privada local. Uma thread dedicada à recepção permite receber mensagens enquanto o usuário digita (tempo real).
\end{enumerate}

\subsection{Exemplo de log de uma conversa}
O trecho abaixo, gerado em teste real, mostra no servidor o handshake e as mensagens. Observe que o que trafega na rede é apenas hexadecimal ilegível:

\begin{lstlisting}[style=logstyle]
[23:26:00] [SERVIDOR] Aguardando cliente na porta 5071...
[23:26:01] [SERVIDOR] Cliente conectado: 127.0.0.1:55096
[23:26:01] [HANDSHAKE] Enviando minha chave publica ao outro no...
[23:26:01] [HANDSHAKE] Chave publica do outro no recebida:
[23:26:01] [HANDSHAKE] e = 65537
[23:26:01] [HANDSHAKE] n = ... (256 hex)
[23:26:02] [HANDSHAKE] Handshake concluido. Digite mensagens (/sair para encerrar).
[23:26:02] [CIFRADO] "ola cliente, tudo bem?" -> cifrado com a chave publica do destinatario
[23:26:02] [REDE ->] Enviando: 9dbd3cc8534b9defb5ceb16b22634af20907fec8215b54b39174774f9437... (256 hex)
[23:26:04] [REDE <-] Cifra recebida: 89208e8af87210cb7a451d6a059a6d00adb874f8ab89ce4d479a108048c9... (256 hex)
[23:26:04] [DECIFRADO] Usando minha chave privada -> "oi servidor! cao"
\end{lstlisting}

Qualquer interceptação do tráfego na porta do chat (por exemplo, com Wireshark) exibirá somente as linhas \texttt{PUBKEY} (públicas por natureza) e \texttt{MSG} com a cifra em hexadecimal, sem nenhum texto legível.

\section{Conclusão}
O trabalho permitiu relacionar a teoria e a prática do RSA: da escolha de primos com Miller-Rabin e do cálculo de $n$, $\phi(n)$, $e$ e $d$, até o uso do algoritmo em uma aplicação de rede em que as mensagens só trafegam cifradas. A implementação é didática: demonstra os princípios, mas sistemas reais exigem chaves de pelo menos 2048 bits, preenchimento OAEP, autenticação de chaves por certificados e, cada vez mais, algoritmos resistentes à computação quântica.

\section*{Referências}
\begin{enumerate}[label={[\arabic*]}, leftmargin=*]
    \item RIVEST, R. L.; SHAMIR, A.; ADLEMAN, L. A method for obtaining digital signatures and public-key cryptosystems. \textbf{Communications of the ACM}, v. 21, n. 2, p. 120-126, 1978.
    \item Material de referência principal indicado pelo professor (link do Google Drive fornecido no enunciado do trabalho).
    \item BOUDOT, F. et al. Comparing the difficulty of factorization and discrete logarithm: a 240-digit experiment. \textbf{CRYPTO 2020} (mesma equipe que anunciou a fatoração do RSA-250, em 2020).
    \item BARKER, E. \textbf{Recommendation for Key Management: Part 1 - General}. NIST Special Publication 800-57 Part 1.
    \item SHOR, P. W. Polynomial-time algorithms for prime factorization and discrete logarithms on a quantum computer. \textbf{SIAM Journal on Computing}, v. 26, n. 5, 1997.
    \item RSA Laboratories. \textbf{PKCS \#1: RSA Cryptography Specifications}; RFC 8017.
    \item BONEH, D. Twenty years of attacks on the RSA cryptosystem. \textbf{Notices of the AMS}, v. 46, n. 2, 1999.
    \item DIFFIE, W.; HELLMAN, M. New directions in cryptography. \textbf{IEEE Transactions on Information Theory}, v. 22, n. 6, 1976.
\end{enumerate}

\end{document}